\section{Analysis and Point-Set Topology}

We provide no proofs here.

\subsection{Real Analysis}


\subsection{Topology}

\subsubsection{Open and Closed Sets}

\begin{lemma}
    Let $Y$ be a subspace of $X$. If $U$ is open in $Y$ and $Y$ is open in $X$, then
    $U$ is open in $X$.
\end{lemma}

\begin{remark}
    Note that if $Y$ is not open in $X$, then there is no reason $U$ should be open in $X$.
    As a counter example, consider $X = \R$ with the standard topology, $Y = [0,1]$ (with the 
    subpace topology), and $U = [0,\frac{1}{2})$. $U = (-\frac{1}{2}, \frac{1}{2}) \cap Y$, 
    so is open in $Y$. However, clearly, $U$ is not open in $X$, since opens in $X$ are unions
    of open intervals.
\end{remark}

\begin{example}
    In the discrete topology on a set $X$, every subset of $X$ is open.
    It follows that every set is closed as well.
\end{example}

\begin{example}
    In any topology space $X$, $\emptyset$ and $X$ are both closed and open.
\end{example}

A nice way to test if a set is closed in the subspace topology of some 
larger set is from the following:

\begin{theorem}
    Let $Y \subseteq X$. $A$ is closed in $Y$ if and only if it equals 
    the intersection of a closed set of $X$ with $Y$.
\end{theorem}

It also follows then, from the lemma earlier, that:
\begin{corollary}
    Let $Y \subseteq X$. If $A$ is closed in $Y$ and $Y$ is closed in $X$, 
    then $A$ is closed in $X$.
\end{corollary}

There was also the notion of a closure, which is the intersection of all closed
sets containing a set. Note that the resulting set is closed because arbitrary 
intersections of closed sets are closed. 

\begin{theorem}
    Let $A \subseteq X$ be (only) a subset.
    \begin{itemize}
        \item $x \in \overline{A}$ if and only if every open set $U$ containing $x$
            intersects $A$.
        \item If the topology of $X$ is given by a basis, then $x \in \overline{A}$ 
        if and only if every basis element containing $x$ intersects $A$.
    \end{itemize}
\end{theorem}

\begin{exercise}
    Let $X = \{n \in \Z \colon n \geq 2\}$ be the topoological space $X$ with the topology 
    generated by sets of the form $U_k := \{n \in X \colon n \mid k\}$, for $k \geq 2$, so that
    open sets of $X$ are the empty set and arbitrary unions of the sets $U_k$. For any $n \in X$,
    determine the closure of the one point set $\{n\}$ in $X$.
\end{exercise}
\begin{solution}
    Note that $m \in \overline{\{n\}}$ if and only if every open neighborhood of $m$ meets
    $\{n\}$. In other words, we mean that if $m \in U_k$, then $\{n\} \cap U_k \neq \emptyset$,
    so $n \in U_k$. That is to say, if $m \mid k$, then we need $n \mid k$. Therefore, $m = ns$.
    Therefore, the closure is of the form $\{sn \colon s \in \Z_{>0}\}$.
\end{solution}

\begin{df}
    If $A$ is a subset of the topological space $X$, then $x \in X$ is called a \emph{limit point}
    of $A$ if every neighborhood of $x$ intersects $A$ in some point other than $x$ itself.
\end{df}

\begin{remark}
    In other words, $x$ is a limit point of $A$ if $x$ is in the closure of $A \setminus \{x\}$.
\end{remark}

\begin{theorem}
    The closure of a set $A$ in a topological space $X$ is precisely the set $A$ union with
    the set of limit points of $A$ in $X$.
\end{theorem}

\subsubsection{Hausdorff}

Recall the notion of Hausdorff:

\begin{df}
    A topological space $X$ is \emph{Hausdorff} if for each pair of distinct 
    points $x,y \in X$, there are neighborhoods of $x,y$, respectively, 
    which are disjoint.
\end{df}

\begin{theorem}
    Every finite point set in a Hausdorff space is closed.
\end{theorem}

\begin{theorem}
    If $X$ is a Hausdorff space, then a sequence of points of $X$ converges 
    to at most one point of $X$.
\end{theorem}

\begin{proposition}
    Products and subspaces of Hausdorff spaces are Hausdorff.
\end{proposition}

\subsection{Miscellanious}